%%%%%%%%%%%%%%%%%%%%%%%%%%%%%%%%%%%%%%%%%
% Masters/Doctoral Thesis 
% LaTeX Template
% Version 1.43 (17/5/14)
%
% This template has been downloaded from:
% http://www.LaTeXTemplates.com
%
% Original authors:
% Steven Gunn 
% http://users.ecs.soton.ac.uk/srg/softwaretools/document/templates/
% and
% Sunil Patel
% http://www.sunilpatel.co.uk/thesis-template/
%
% License:
% CC BY-NC-SA 3.0 (http://creativecommons.org/licenses/by-nc-sa/3.0/)
%
% Note:
% Make sure to edit document variables in the Thesis.cls file
%
%%%%%%%%%%%%%%%%%%%%%%%%%%%%%%%%%%%%%%%%%

%----------------------------------------------------------------------------------------
%	PACKAGES AND OTHER DOCUMENT CONFIGURATIONS
%----------------------------------------------------------------------------------------

\documentclass[11pt, oneside]{Thesis} % The default font size and one-sided printing (no margin offsets)

\graphicspath{{Pictures/}} % Specifies the directory where pictures are stored

\usepackage[square, numbers, comma, sort&compress]{natbib} % Use the natbib reference package - read up on this to edit the reference style; if you want text (e.g. Smith et al., 2012) for the in-text references (instead of numbers), remove 'numbers' 
\usepackage{dirtytalk, cleveref,listings, color, graphicx,pxfonts,amsmath}
\hypersetup{urlcolor=black, colorlinks=true} % Colors hyperlinks in blue - change to black if annoying
\title{\ttitle} % Defines the thesis title - don't touch this

\definecolor{dkgreen}{rgb}{0,0.6,0}
\definecolor{gray}{rgb}{0.5,0.5,0.5}
\definecolor{mauve}{rgb}{0.58,0,0.82}
\definecolor{gray}{rgb}{0.4,0.4,0.4}
\definecolor{darkblue}{rgb}{0.0,0.0,0.6}
\definecolor{lightblue}{rgb}{0.0,0.0,0.9}
\definecolor{cyan}{rgb}{0.0,0.6,0.6}
\definecolor{darkred}{rgb}{0.6,0.0,0.0}

\lstset{
  basicstyle=\ttfamily\scriptsize,
  columns=fullflexible,
  showstringspaces=false,
  keywordstyle=\color{black}\bfseries,
  numbers=left,                   % where to put the line-numbers
  numberstyle=\tiny\color{gray},  % the style that is used for the line-numbers
  stepnumber=1,
  numbersep=5pt,                  % how far the line-numbers are from the code
  backgroundcolor=\color{white},      % choose the background color. You must add \usepackage{color}
  showspaces=false,               % show spaces adding particular underscores
  showstringspaces=false,         % underline spaces within strings
  showtabs=false,                 % show tabs within strings adding particular underscores
  frame=none,                   % adds a frame around the code
  rulecolor=\color{black},        % if not set, the frame-color may be changed on line-breaks within not-black text (e.g. commens (green here))
  tabsize=2,                      % sets default tabsize to 2 spaces
  captionpos=b,                   % sets the caption-position to bottom
  breaklines=true,                % sets automatic line breaking
  breakatwhitespace=false,        % sets if automatic breaks should only happen at whitespace
  title=\lstname,                   % show the filename of files included with \lstinputlisting;
                                  % also try caption instead of title
  commentstyle=\color{gray}\upshape
}

\lstdefinelanguage{SPARQL}{
	morekeywords={PREFIX, SELECT, WHERE, FILTER, OPTIONAL, DISTINCT}
}

\lstdefinelanguage{XML}
{
  morestring=[s][\color{mauve}]{"}{"},
  morestring=[s][\color{black}]{>}{<},
  morecomment=[s]{<?}{?>},
  morecomment=[s][\color{dkgreen}]{!--}{--},
  stringstyle=\color{black},
  identifierstyle=\color{lightblue},
  keywordstyle=\color{red},
  % morekeywords={xmlns,xsi,noNamespaceSchemaLocation,type,id,x,y,source,target,version,tool,transRef,roleRef,objective,eventually}% list your attributes here
}

\lstdefinelanguage{json}{}

\begin{document}

\frontmatter % Use roman page numbering style (i, ii, iii, iv...) for the pre-content pages

\setstretch{1.3} % Line spacing of 1.3

% Define the page headers using the FancyHdr package and set up for one-sided printing
\fancyhead{} % Clears all page headers and footers
\rhead{\thepage} % Sets the right side header to show the page number
\lhead{} % Clears the left side page header

\pagestyle{fancy} % Finally, use the "fancy" page style to implement the FancyHdr headers

\newcommand{\HRule}{\rule{\linewidth}{0.5mm}} % New command to make the lines in the title page
\renewcommand{\lstlistingname}{Code}

% PDF meta-data
\hypersetup{pdftitle={\ttitle}}
\hypersetup{pdfsubject=\subjectname}
\hypersetup{pdfauthor=\authornames}
\hypersetup{pdfkeywords=\keywordnames}

%----------------------------------------------------------------------------------------
%	TITLE PAGE
%----------------------------------------------------------------------------------------

\begin{titlepage}
\begin{center}

\textsc{\LARGE \univname}\\[1.5cm] % University name
\textsc{\Large Master Thesis}\\[0.5cm] % Thesis type

\HRule \\[0.4cm] % Horizontal line
{\huge \bfseries \ttitle}\\[0.4cm] % Thesis title
\HRule \\[1.5cm] % Horizontal line
 
\begin{minipage}{0.4\textwidth}
\begin{flushleft} \large
\emph{Author:}\\
\authornames % Author name - remove the \href bracket to remove the link
\end{flushleft}
\end{minipage}
\begin{minipage}{0.4\textwidth}
\begin{flushright} \large
\emph{Supervisors:} \\
\supname % Supervisor name - remove the \href bracket to remove the link  
\end{flushright}
\end{minipage}\\[3cm]
 
\large \textit{This thesis submitted in fulfilment of the requirements\\ for the degree of \degreename}\\[0.3cm] % University requirement text
\textit{in the}\\[0.4cm]
\groupname\\\deptname\\[2cm] % Research group name and department name
 
{\large \today}\\[4cm] % Date
%\includegraphics{Logo} % University/department logo - uncomment to place it
 
\vfill
\end{center}

\end{titlepage}

%----------------------------------------------------------------------------------------
% DECLARATION PAGE
% Your institution may give you a different text to place here
%----------------------------------------------------------------------------------------

\Declaration{

\addtocontents{toc}{\vspace{1em}} % Add a gap in the Contents, for aesthetics

I, \authornames, declare that this thesis titled, '\ttitle' and the work presented in it are my own. I confirm that:

\begin{itemize} 
\item[\tiny{$\blacksquare$}] This work was done wholly or mainly while in candidature for a research degree at this University.
\item[\tiny{$\blacksquare$}] Where any part of this thesis has previously been submitted for a degree or any other qualification at this University or any other institution, this has been clearly stated.
\item[\tiny{$\blacksquare$}] Where I have consulted the published work of others, this is always clearly attributed.
\item[\tiny{$\blacksquare$}] Where I have quoted from the work of others, the source is always given. With the exception of such quotations, this thesis is entirely my own work.
\item[\tiny{$\blacksquare$}] I have acknowledged all main sources of help.
\item[\tiny{$\blacksquare$}] Where the thesis is based on work done by myself jointly with others, I have made clear exactly what was done by others and what I have contributed myself.\\
\end{itemize}
 
Signed:\\
\rule[1em]{25em}{0.5pt} % This prints a line for the signature
 
Date:\\
\rule[1em]{25em}{0.5pt} % This prints a line to write the date
}

\clearpage % Start a new page

%----------------------------------------------------------------------------------------
%	QUOTATION PAGE
%----------------------------------------------------------------------------------------

\pagestyle{empty} % No headers or footers for the following pages

\null\vfill % Add some space to move the quote down the page a bit

\textit{``I have a dream for the Web [in which computers] become capable of analyzing all the data on the Web--the content, links, and transactions between people and computers. A ``Semantic Web", which should make this possible, has yet to emerge, but when it does, the day-to-day mechanisms of trade, bureaucracy and our daily lives will be handled by machines talking to machines."}

\begin{flushright}
Sir Tim Berners-Lee
\end{flushright}

\vfill\vfill\vfill\vfill\vfill\vfill\null % Add some space at the bottom to position the quote just right

\clearpage % Start a new page

%----------------------------------------------------------------------------------------
%	ACKNOWLEDGEMENTS
%----------------------------------------------------------------------------------------

\setstretch{1.3} % Reset the line-spacing to 1.3 for body text (if it has changed)

\acknowledgements{\addtocontents{toc}{\vspace{1em}} % Add a gap in the Contents, for aesthetics


This master thesis was done in the context of the double master degree Informatique - Information und Kommunikation (IFIK), which brings together two Master programs : a degree in computer engineering at the National Institute of Applied Sciences in Lyon (INSA Lyon) and a Master in Informatik (Schwerpunkt : Information und Kommunikationssysteme) at the University of Passau.\\

First, I would like to thanks the initiators of the double Master IFIK for their valuable efforts and commitment to the german-french collaboration : Prof. Dr. Harald Kosch for the german side and Prof. Dr. Lionel Brunie for the french one.\\

An unmeasurable amout of thanks goes to my advisors, Dr. David Coquil and Dr. Elöd Egyed-Zsigmond for their precious help and advices during this thesis preparation.\\

I would also thanks my colleagues and friends for their support as well as my five german roommates who made my time in Passau such enjoyable.\\

Finally, I thanks the Dropkick Murphys for their energizing music and the worldwide coffee producers for this magical beverage they provide us.
}
\clearpage % Start a new page

%----------------------------------------------------------------------------------------
%	ABSTRACT PAGE
%----------------------------------------------------------------------------------------

\addtotoc{Abstract} % Add the "Abstract" page entry to the Contents

\abstract{\addtocontents{toc}{\vspace{1em}} % Add a gap in the Contents, for aesthetics

Image is a popular medium nowadays : it is easy to capture, can be really light on your electronic device and speaks to everyone without distinction of language. Lot of companies use image data every day and so need a good way to retrieve them. For a human-being, the more intuitive way to describe what he's looking for is by the use of his own words, not a distribution of colors or any low-level feature. Therefore, it is interesting to add annotations (also called tags) to images in order to describe them in an natural language way. In this work, we propose a semantic enrichment prototype based on the semantic web and graph models. Five experiments are presented and their results are discussed. Eventually we expose some future work which should be interesting.
\keywords{image annotation, semantic enrichment, semantic web, graph model}
}

\clearpage % Start a new page

%----------------------------------------------------------------------------------------
%	LIST OF CONTENTS/FIGURES/TABLES PAGES
%----------------------------------------------------------------------------------------

\pagestyle{fancy} % The page style headers have been "empty" all this time, now use the "fancy" headers as defined before to bring them back

\lhead{\emph{Contents}} % Set the left side page header to "Contents"
\tableofcontents % Write out the Table of Contents

\lhead{\emph{List of Figures}} % Set the left side page header to "List of Figures"
\listoffigures % Write out the List of Figures

\lhead{\emph{List of Tables}} % Set the left side page header to "List of Tables"
\listoftables % Write out the List of Tables

%----------------------------------------------------------------------------------------
%	ABBREVIATIONS
%----------------------------------------------------------------------------------------

\clearpage % Start a new page

\setstretch{1.5} % Set the line spacing to 1.5, this makes the following tables easier to read

\lhead{\emph{Abbreviations}} % Set the left side page header to "Abbreviations"
\listofsymbols{ll} % Include a list of Abbreviations (a table of two columns)
{
\textbf{LCS} & \textbf{L}east \textbf{C}ommon \textbf{S}ubsumer \\
\textbf{RDF} & \textbf{R}esource \textbf{D}escription \textbf{F}ramework \\
\textbf{SP} & \textbf{S}hortest \textbf{P}ath \\
\textbf{SPARQL} & \textbf{S}PARQL \textbf{P}rotocol \textbf{a}nd \textbf{R}DF \textbf{Q}uery \textbf{L}anguage \\
\textbf{URI} & \textbf{U}niversal \textbf{R}esource \textbf{I}dentifier \\
\textbf{WP} & \textbf{W}u-\textbf{P}almer measure \\
\textbf{WP++} & \textbf{W}u-\textbf{P}almer evolved measure \\
\textbf{XML} & e\textbf{X}tensible \textbf{M}arkup \textbf{L}anguage \\
}

%----------------------------------------------------------------------------------------
%	THESIS CONTENT - CHAPTERS
%----------------------------------------------------------------------------------------

\mainmatter % Begin numeric (1,2,3...) page numbering

\pagestyle{fancy} % Return the page headers back to the "fancy" style

% Include the chapters of the thesis as separate files from the Chapters folder
% Uncomment the lines as you write the chapters

% % Chapter 1

\chapter{Chapter Title Here} % Main chapter title

\label{Chapter1} % For referencing the chapter elsewhere, use \ref{Chapter1} 

\lhead{Chapter 1. \emph{Chapter Title Here}} % This is for the header on each page - perhaps a shortened title

%----------------------------------------------------------------------------------------

\section{Welcome and Thank You}
Welcome to this \LaTeX{} Thesis Template, a beautiful and easy to use template for writing a thesis using the \LaTeX{} typesetting system.

If you are writing a thesis (or will be in the future) and its subject is technical or mathematical (though it doesn't have to be), then creating it in \LaTeX{} is highly recommended as a way to make sure you can just get down to the essential writing without having to worry over formatting or wasting time arguing with your word processor.

\LaTeX{} is easily able to professionally typeset documents that run to hundreds or thousands of pages long. With simple mark-up commands, it automatically sets out the table of contents, margins, page headers and footers and keeps the formatting consistent and beautiful. One of its main strengths is the way it can easily typeset mathematics, even \emph{heavy} mathematics. Even if those equations are the most horribly twisted and most difficult mathematical problems that can only be solved on a super-computer, you can at least count on \LaTeX{} to make them look stunning.

%----------------------------------------------------------------------------------------

\section{Learning \LaTeX{}}

\LaTeX{} is not a WYSIWYG (What You See is What You Get) program, unlike word processors such as Microsoft Word or Apple's Pages. Instead, a document written for \LaTeX{} is actually a simple, plain text file that contains \emph{no formatting}. You tell \LaTeX{} how you want the formatting in the finished document by writing in simple commands amongst the text, for example, if I want to use \textit{italic text for emphasis}, I write the `$\backslash$\texttt{textit}\{\}' command and put the text I want in italics in between the curly braces. This means that \LaTeX{} is a ``mark-up'' language, very much like HTML.

\subsection{A (not so short) Introduction to \LaTeX{}}

If you are new to \LaTeX{}, there is a very good eBook -- freely available online as a PDF file -- called, ``The Not So Short Introduction to \LaTeX{}''. The book's title is typically shortened to just ``lshort''. You can download the latest version (as it is occasionally updated) from here:\\
\href{http://www.ctan.org/tex-archive/info/lshort/english/lshort.pdf}{\texttt{http://www.ctan.org/tex-archive/info/lshort/english/lshort.pdf}}

It is also available in several other languages. Find yours from the list on this page:\\
\href{http://www.ctan.org/tex-archive/info/lshort/}{\texttt{http://www.ctan.org/tex-archive/info/lshort/}}

It is recommended to take a little time out to learn how to use \LaTeX{} by creating several, small `test' documents. Making the effort now means you're not stuck learning the system when what you \emph{really} need to be doing is writing your thesis.

\subsection{A Short Math Guide for \LaTeX{}}

If you are writing a technical or mathematical thesis, then you may want to read the document by the AMS (American Mathematical Society) called, ``A Short Math Guide for \LaTeX{}''. It can be found online here:\\
\href{http://www.ams.org/tex/amslatex.html}{\texttt{http://www.ams.org/tex/amslatex.html}}\\
under the ``Additional Documentation'' section towards the bottom of the page.

\subsection{Common \LaTeX{} Math Symbols}
There are a multitude of mathematical symbols available for \LaTeX{} and it would take a great effort to learn the commands for them all. The most common ones you are likely to use are shown on this page:\\
\href{http://www.sunilpatel.co.uk/latexsymbols.html}{\texttt{http://www.sunilpatel.co.uk/latexsymbols.html}}

You can use this page as a reference or crib sheet, the symbols are rendered as large, high quality images so you can quickly find the \LaTeX{} command for the symbol you need.

\subsection{\LaTeX{} on a Mac}
 
The \LaTeX{} package is available for many systems including Windows, Linux and Mac OS X. The package for OS X is called MacTeX and it contains all the applications you need -- bundled together and pre-customised -- for a fully working \LaTeX{} environment and workflow.
 
MacTeX includes a dedicated \LaTeX{} IDE (Integrated Development Environment) called ``TeXShop'' for writing your `\texttt{.tex}' files and ``BibDesk'': a program to manage your references and create your bibliography section just as easily as managing songs and creating playlists in iTunes.

%----------------------------------------------------------------------------------------

\section{Getting Started with this Template}

If you are familiar with \LaTeX{}, then you can familiarise yourself with the contents of the Zip file and the directory structure and then place your own information into the `\texttt{Thesis.cls}' file. Section \ref{FillingFile} on page \pageref{FillingFile} tells you how to do this. Make sure you read section \ref{ThesisConventions} about thesis conventions to get the most out of this template and then get started with the `\texttt{Thesis.tex}' file straightaway.

If you are new to \LaTeX{} it is recommended that you carry on reading through the rest of the information in this document.

\subsection{About this Template}

This \LaTeX{} Thesis Template is originally based and created around a \LaTeX{} style file created by Steve R.\ Gunn from the University of Southampton (UK), department of Electronics and Computer Science. You can find his original thesis style file at his site, here:\\
\href{http://www.ecs.soton.ac.uk/~srg/softwaretools/document/templates/}{\texttt{http://www.ecs.soton.ac.uk/$\sim$srg/softwaretools/document/templates/}}

My thesis originally used the `\texttt{ecsthesis.cls}' from his list of styles. However, I knew \LaTeX{} could still format better. To get the look I wanted, I modified his style and also created a skeleton framework and folder structure to place the thesis files in.

This Thesis Template consists of that modified style, the framework and the folder structure. All the work that has gone into the preparation and groundwork means that all you have to bother about is the writing.

Before you begin using this template you should ensure that its style complies with the thesis style guidelines imposed by your institution. In most cases this template style and layout will be suitable. If it is not, it may only require a small change to bring the template in line with your institution's recommendations.

%----------------------------------------------------------------------------------------

\section{What this Template Includes}

\subsection{Folders}

This template comes as a single Zip file that expands out to many files and folders. The folder names are mostly self-explanatory:

\textbf{Appendices} -- this is the folder where you put the appendices. Each appendix should go into its own separate `\texttt{.tex}' file. A template is included in the directory.

\textbf{Chapters} -- this is the folder where you put the thesis chapters. A thesis usually has about seven chapters, though there is no hard rule on this. Each chapter should go in its own separate `\texttt{.tex}' file and they usually are split as:
\begin{itemize}
\item Chapter 1: Introduction to the thesis topic
\item Chapter 2: Background information and theory
\item Chapter 3: (Laboratory) experimental setup
\item Chapter 4: Details of experiment 1
\item Chapter 5: Details of experiment 2
\item Chapter 6: Discussion of the experimental results
\item Chapter 7: Conclusion and future directions
\end{itemize}
This chapter layout is specialised for the experimental sciences.

\textbf{Figures} -- this folder contains all figures for the thesis. These are the final images that will go into the thesis document.

\textbf{Primitives} -- this is the folder that contains scraps, particularly because one final image in the `Figures' folder may be made from many separate images and photos, these source images go here. This keeps the intermediate files separate from the final thesis figures.

\subsection{Files}

Included are also several files, most of them are plain text and you can see their contents in a text editor. Luckily, many of them are auxiliary files created by \LaTeX{} or BibTeX and which you don't need to bother about:

\textbf{Bibliography.bib} -- this is an important file that contains all the bibliographic information and references that you will be citing in the thesis for use with BibTeX. You can write it manually, but there are reference manager programs available that will create and manage it for you. Bibliographies in \LaTeX{} are a large subject and you may need to read about BibTeX before starting with this.

\textbf{Thesis.cls} -- this is an important file. It is the style file that tells \LaTeX{} how to format the thesis. You will also need to open this file in a text editor and fill in your own information (such as name, department, institution). Luckily, this is not too difficult and is explained in section \ref{FillingFile} on page \pageref{FillingFile}.

\textbf{Thesis.pdf} -- this is your beautifully typeset thesis (in the PDF file format) created by \LaTeX{}.

\textbf{Thesis.tex} -- this is an important file. This is the file that you tell \LaTeX{} to compile to produce your thesis as a PDF file. It contains the framework and constructs that tell \LaTeX{} how to layout the thesis. It is heavily commented so you can read exactly what each line of code does and why it is there. After you put your own information into the `\texttt{Thesis.cls}' file, go to this file and begin filling it in -- you have now started your thesis!

\textbf{vector.sty} -- this is a \LaTeX{} package, it tells \LaTeX{} how to typeset mathematical vectors. Using this package is very easy and you can read the documentation on the site (you just need to look at the `\texttt{vector.pdf}' file):\\
\href{http://www.ctan.org/tex-archive/macros/latex/contrib/vector/}{\texttt{http://www.ctan.org/tex-archive/macros/latex/contrib/vector/}}

\textbf{lstpatch.sty} -- this is a \LaTeX{} package required by this LaTeX template and is included as not all \TeX{} distributions have it installed by default. You do not need to modify this file.

Files that are \emph{not} included, but are created by \LaTeX{} as auxiliary files include:

\textbf{Thesis.aux} -- this is an auxiliary file generated by \LaTeX{}, if it is deleted \LaTeX{} simply regenerates it when you run the main `\texttt{.tex}' file.

\textbf{Thesis.bbl} -- this is an auxiliary file generated by BibTeX, if it is deleted, BibTeX simply regenerates it when you run the main tex file. Whereas the `\texttt{.bib}' file contains all the references you have, this `\texttt{.bbl}' file contains the references you have actually cited in the thesis and is used to build the bibliography section of the thesis.

\textbf{Thesis.blg} -- this is an auxiliary file generated by BibTeX, if it is deleted BibTeX simply regenerates it when you run the main `\texttt{.tex}' file.

\textbf{Thesis.lof} -- this is an auxiliary file generated by \LaTeX{}, if it is deleted \LaTeX{} simply regenerates it when you run the main `\texttt{.tex}' file. It tells \LaTeX{} how to build the `List of Figures' section.

\textbf{Thesis.log} -- this is an auxiliary file generated by \LaTeX{}, if it is deleted \LaTeX{} simply regenerates it when you run the main `\texttt{.tex}' file. It contains messages from \LaTeX{}, if you receive errors and warnings from \LaTeX{}, they will be in this `\texttt{.log}' file.

\textbf{Thesis.lot} -- this is an auxiliary file generated by \LaTeX{}, if it is deleted \LaTeX{} simply regenerates it when you run the main `\texttt{.tex}' file. It tells \LaTeX{} how to build the `List of Tables' section.

\textbf{Thesis.out} -- this is an auxiliary file generated by \LaTeX{}, if it is deleted \LaTeX{} simply regenerates it when you run the main `\texttt{.tex}' file.


So from this long list, only the files with the `\texttt{.sty}', `\texttt{.bib}', `\texttt{.cls}' and `\texttt{.tex}' extensions are the most important ones. The other auxiliary files can be ignored or deleted as \LaTeX{} and BibTeX will regenerate them.

%----------------------------------------------------------------------------------------

\section{Filling in the `\texttt{Thesis.cls}' File}\label{FillingFile}

You will need to personalise the thesis template and make it your own by filling in your own information. This is done by editing the `\texttt{Thesis.cls}' file in a text editor.

Open the file and scroll down, past all the `$\backslash$\texttt{newcommand}\ldots' items until you see the entries for `\texttt{University Name}', `\texttt{Department Name}', etc\ldots.

Fill out the information about your group and institution and ensure you keep to block capitals where it asks you to. You can also insert web links, if you do, make sure you use the full URL, including the `\texttt{http://}' for this.

The last item you should need to fill in is the Faculty Name (in block capitals). When you have done this, save the file and recompile `\texttt{Thesis.tex}'. All the information you filled in should now be in the PDF, complete with web links. You can now begin your thesis proper!

%----------------------------------------------------------------------------------------

\section{The `\texttt{Thesis.tex}' File Explained}

The \texttt{Thesis.tex} file contains the structure of the thesis. There are plenty of written comments that explain what pages, sections and formatting the \LaTeX{} code is creating. Initially there seems to be a lot of \LaTeX{} code, but this is all formatting, and it has all been taken care of so you don't have to do it.

Begin by checking that your information on the title page is correct. For the thesis declaration, your institution may insist on something different than the text given. If this is the case, just replace what you see with what is required.

Then comes a page which contains a funny quote. You can put your own, or quote your favourite scientist, author, person, etc\ldots Make sure to put the name of the person who you took the quote from.

Next comes the acknowledgements. On this page, write about all the people who you wish to thank (not forgetting parents, partners and your advisor/supervisor).

The contents pages, list of figures and tables are all taken care of for you and do not need to be manually created or edited. The next set of pages are optional and can be deleted since they are for a more technical thesis: insert a list of abbreviations you have used in the thesis, then a list of the physical constants and numbers you refer to and finally, a list of mathematical symbols used in any formulae. Making the effort to fill these tables means the reader has a one-stop place to refer to instead of searching the internet and references to try and find out what you meant by certain abbreviations or symbols.

The list of symbols is split into the Roman and Greek alphabets. Whereas the abbreviations and symbols ought to be listed in alphabetical order (and this is \emph{not} done automatically for you) the list of physical constants should be grouped into similar themes.

The next page contains a one line dedication. Who will you dedicate your thesis to?

Finally, there is the section where the chapters are included. Uncomment the lines (delete the `\texttt{\%}' character) as you write the chapters. Each chapter should be written in its own file and put into the `Chapters' folder and named `\texttt{Chapter1}', `\texttt{Chapter2}, etc\ldots Similarly for the appendices, uncomment the lines as you need them. Each appendix should go into its own file and placed in the `Appendices' folder.

After the preamble, chapters and appendices finally comes the bibliography. The bibliography style (called `\texttt{unsrtnat}') is used for the bibliography and is a fully featured style that will even include links to where the referenced paper can be found online. Do not under estimate how grateful you reader will be to find that a reference to a paper is just a click away. Of course, this relies on you putting the URL information into the BibTeX file in the first place.

%----------------------------------------------------------------------------------------

\section{Thesis Features and Conventions}\label{ThesisConventions}

To get the best out of this template, there are a few conventions that you may want to follow.

One of the most important (and most difficult) things to keep track of in such a long document as a thesis is consistency. Using certain conventions and ways of doing things (such as using a Todo list) makes the job easier. Of course, all of these are optional and you can adopt your own method.

\subsection{Printing Format}

This thesis template is designed for single sided printing as most theses are printed and bound this way. This means that the left margin is always wider than the right (for binding). Four out of five people will now judge the margins by eye and think, ``I never 
noticed that before.''.

The headers for the pages contain the page number on the right side (so it is easy to flick through to the page you want) and the chapter name on the left side.

The text is set to 11 point and a line spacing of 1.3. Generally, it is much more readable to have a smaller text size and wider gap between the lines than it is to have a larger text size and smaller gap. Again, you can tune the text size and spacing should you want or need to. The text size can be set in the options for the `$\backslash$\texttt{documentclass}' command at the top of the `\texttt{Thesis.tex}' file and the spacing can be changed by setting a different value in the `$\backslash$\texttt{setstretch}' commands (scattered throughout the `\texttt{Thesis.tex}' file).

\subsection{Using US Letter Paper}

The paper size used in the template is A4, which is a common -- if not standard -- size in Europe. If you are using this thesis template elsewhere and particularly in the United States, then you may have to change the A4 paper size to the US Letter size. Unfortunately, this is not as simple as replacing instances of `\texttt{a4paper}' with `\texttt{letterpaper}'.

This is because the final PDF file is created directly from the \LaTeX{} source using a program called `\texttt{pdfTeX}' and in certain conditions, paper size commands are ignored and all documents are created with the paper size set to the size stated in the configuration file for pdfTeX (called `\texttt{pdftex.cfg}').

What needs to be done is to change the paper size in the configuration file for \texttt{pdfTeX} to reflect the letter size. There is an excellent tutorial on how to do this here: \\
\href{http://www.physics.wm.edu/~norman/latexhints/pdf_papersize.html}{\texttt{http://www.physics.wm.edu/$\sim$norman/latexhints/pdf\_papersize.html}}

It may be sufficient just to replace the dimensions of the A4 paper size with the US Letter size in the \texttt{pdftex.cfg} file. Due to the differences in the paper size, the resulting margins may be different to what you like or require (as it is common for Institutions to dictate certain margin sizes). If this is the case, then the margin sizes can be tweaked by opening up the \texttt{Thesis.cls} file and searching for the line beginning with, `$\backslash$\texttt{setmarginsrb}' (not very far down from the top), there you will see the margins specified. Simply change those values to what you need (or what looks good) and save. Now your document should be set up for US Letter paper size with suitable margins.

\subsection{References}

The `\texttt{natbib}' package is used to format the bibliography and inserts references such as this one \citep{Reference3}. The options used in the `\texttt{Thesis.tex}' file mean that the references are listed in numerical order as they appear in the text. Multiple references are rearranged in numerical order (e.g. \citep{Reference2, Reference1}) and multiple, sequential references become reformatted to a reference range (e.g. \citep{Reference2, Reference1, Reference3}). This is done automatically for you. To see how you use references, have a look at the `\texttt{Chapter1.tex}' source file. Many reference managers allow you to simply drag the reference into the document as you type.

Scientific references should come \emph{before} the punctuation mark if there is one (such as a comma or period). The same goes for footnotes\footnote{Such as this footnote, here down at the bottom of the page.}. You can change this but the most important thing is to keep the convention consistent throughout the thesis. Footnotes themselves should be full, descriptive sentences (beginning with a capital letter and ending with a full stop).

To see how \LaTeX{} typesets the bibliography, have a look at the very end of this document (or just click on the reference number links).

\subsection{Figures}

There will hopefully be many figures in your thesis (that should be placed in the `Figures' folder). The way to insert figures into your thesis is to use a code template like this:
\begin{verbatim}
\begin{figure}[htbp]
  \centering
    \includegraphics{Figures/Electron.pdf}
    \rule{35em}{0.5pt}
  \caption[An Electron]{An electron (artist's impression).}
  \label{fig:Electron}
\end{figure}
\end{verbatim}
Also look in the source file. Putting this code into the source file produces the picture of the electron that you can see in the figure below.

\begin{figure}[htbp]
	\centering
		\includegraphics{Figures/Electron.pdf}
		\rule{35em}{0.5pt}
	\caption[An Electron]{An electron (artist's impression).}
	\label{fig:Electron}
\end{figure}

Sometimes figures don't always appear where you write them in the source. The placement depends on how much space there is on the page for the figure. Sometimes there is not enough room to fit a figure directly where it should go (in relation to the text) and so \LaTeX{} puts it at the top of the next page. Positioning figures is the job of \LaTeX{} and so you should only worry about making them look good!

Figures usually should have labels just in case you need to refer to them (such as in Figure \ref{fig:Electron}). The `$\backslash$\texttt{caption}' command contains two parts, the first part, inside the square brackets is the title that will appear in the `List of Figures', and so should be short. The second part in the curly brackets should contain the longer and more descriptive caption text.

The `$\backslash$\texttt{rule}' command is optional and simply puts an aesthetic horizontal line below the image. If you do this for one image, do it for all of them.

The \LaTeX{} Thesis Template is able to use figures that are either in the PDF or JPEG file format.

\subsection{Typesetting mathematics}

If your thesis is going to contain heavy mathematical content, be sure that \LaTeX{} will make it look beautiful, even though it won't be able to solve the equations for you.

The ``Not So Short Introduction to \LaTeX{}'' (available \href{http://www.ctan.org/tex-archive/info/lshort/english/lshort.pdf}{here}) should tell you everything you need to know for most cases of typesetting mathematics. If you need more information, a much more thorough mathematical guide is available from the AMS called, ``A Short Math Guide to \LaTeX{}'' and can be downloaded from:\\
\href{ftp://ftp.ams.org/pub/tex/doc/amsmath/short-math-guide.pdf}{\texttt{ftp://ftp.ams.org/pub/tex/doc/amsmath/short-math-guide.pdf}}

There are many different \LaTeX{} symbols to remember, luckily you can find the most common symbols \href{http://www.sunilpatel.co.uk/latexsymbols.html}{here}. You can use the web page as a quick reference or crib sheet and because the symbols are grouped and rendered as high quality images (each with a downloadable PDF), finding the symbol you need is quick and easy.

You can write an equation, which is automatically given an equation number by \LaTeX{} like this:
\begin{verbatim}
\begin{equation}
E = mc^{2}
  \label{eqn:Einstein}
\end{equation}
\end{verbatim}

This will produce Einstein's famous energy-matter equivalence equation:
\begin{equation}
E = mc^{2}
\label{eqn:Einstein}
\end{equation}

All equations you write (which are not in the middle of paragraph text) are automatically given equation numbers by \LaTeX{}. If you don't want a particular equation numbered, just put the command, `$\backslash$\texttt{nonumber}' immediately after the equation.

%----------------------------------------------------------------------------------------

\section{Sectioning and Subsectioning}

You should break your thesis up into nice, bite-sized sections and subsections. \LaTeX{} automatically builds a table of Contents by looking at all the `$\backslash$\texttt{chapter}$\{\}$', `$\backslash$\texttt{section}$\{\}$' and `$\backslash$\texttt{subsection}$\{\}$' commands you write in the source.

The table of Contents should only list the sections to three (3) levels. A `$\backslash$\texttt{chapter}$\{\}$' is level one (1). A `$\backslash$\texttt{section}$\{\}$' is level two (2) and so a `$\backslash$\texttt{subsection}$\{\}$' is level three (3). In your thesis it is likely that you will even use a `$\backslash$\texttt{subsubsection}$\{\}$', which is level four (4). Adding all these will create an unnecessarily cluttered table of Contents and so you should use the `$\backslash$\texttt{subsubsection$^{*}\{\}$}' command instead (note the asterisk). The asterisk ($^{*}$) tells \LaTeX{} to omit listing the subsubsection in the Contents, keeping it clean and tidy.

%----------------------------------------------------------------------------------------

\section{In Closing}

You have reached the end of this mini-guide. You can now rename or overwrite this pdf file and begin writing your own `\texttt{Chapter1.tex}' and the rest of your thesis. The easy work of setting up the structure and framework has been taken care of for you. It's now your job to fill it out!

Good luck and have lots of fun!

\begin{flushright}
Guide written by ---\\
Sunil Patel: \href{http://www.sunilpatel.co.uk}{www.sunilpatel.co.uk}
\end{flushright}

%!TEX root = ../thesis.tex

\chapter{Introduction} % Main chapter title

\label{Chapter1} % Change X to a consecutive number; for referencing this chapter elsewhere, use \ref{ChapterX}

\lhead{Chapter 1 \emph{Introduction}} % Change X to a consecutive number; this is for the header on each page - perhaps a shortened title

%----------------------------------------------------------------------------------------
%	SECTION 1
%----------------------------------------------------------------------------------------

\section{Background}

Image is a popular medium nowadays : it is easy to capture, can be really light on your computer and speaks to everyone without distinction of language.\\

In the all days life, people share their pictures on social networks in less than a blink of eye. In average, 70M of pictures are posted on Instagram each day and the users hit the \say{Like} button 2.5B times\footnote{Stats from : https://instagram.com/press/}. Other services like Picasa or Flickr exists but aren't as used as Instagram which is the favorite in the eyes of the teen public.\\
Companies also produce a lot of media data. Industry companies need their products' pictures, marketing and advertising studios use a lot of images in order to create new stuff for their client, \dots But the most consumer of media data are obviously mass media themselves : Newspapers, TV shows, news broadcasts are dealing with pictures at every moment of their day.\\

This huge production and consumption of images implies the need of an efficient way to store and search for the relevant one when the time comes. The best illustration to this need is to think of the nice but long moments one had with its relatives searching for the good picture of the new-born nephew in the family pictures album.\\
Since an image itself doesn't have a natural plain-text representation the best way to describe it is to add meta-data (data about the data) such as its date of creation, its dimensions or, and this is what this thesis is about, some tags.\\

There are a lot of ways if one wants to annotate pictures. We can do it manually, using our own words (like \say{Dad}, \say{Home} \dots), we can also analyze the raw picture, its pixel representation and compare some metrics (like the color histogram) to sample images in order to detect known concepts. Moreover, if the image already possesses annotations, we can enrich it semantically. \\

This field is so wide that it is impossible to speak about all the possibilities and technologies. In this study, we will focus on the last point and investigate the automation of the semantic enrichment. We will study the resources at our disposal and propose a solution keeping in mind the facts cited previously.\\

In the following section, we will present and discuss an application scenario to illustrate the motivation behind this thesis. 

%----------------------------------------------------------------------------------------
%	SECTION 2
%----------------------------------------------------------------------------------------

\section{Motivation}

NewsTV is a famous TV news channel which runs 24/7 and only speaks about the current news. It has lot of reporters worldwide, covering the important local news and sending their production to the main site in Paris, France.\\

The employees often need to consult older coverages in order to explain the context of the news, to make the necrology of a famous actor who recently died or to re-use common shots. Therefor, they need to query the central multimedia database management system using keywords they are familiar with like \say{Elections, France, 2007, José Bové}. But sometimes, their research aren't so specific and they are looking for more generic pictures, let say \say{Land, Tree, Animal}.\\

The first kind of keywords had been tagged by the former reporter who produced the coverage but he logically didn't think to add generic terms. NewsTV needs something to do it automatically when a picture, or any media, is first added to its system which a couple of initial tags.\\

Details about which kind of technology can be 	used to achieve this automatic tagging will come in the following sections. To summarize, the goal of this thesis is to propose a running prototype and evaluate different methods of tagging. The questions that we will try to answer during this study are described in the following section.

\section{Thesis Objectives}

Sed ullamcorper quam eu nisl interdum at interdum enim egestas. Aliquam placerat justo sed lectus lobortis ut porta nisl porttitor. Vestibulum mi dolor, lacinia molestie gravida at, tempus vitae ligula. Donec eget quam sapien, in viverra eros. Donec pellentesque justo a massa fringilla non vestibulum metus vestibulum. Vestibulum in orci quis felis tempor lacinia. Vivamus ornare ultrices facilisis. Ut hendrerit volutpat vulputate. Morbi condimentum venenatis augue, id porta ipsum vulputate in. Curabitur luctus tempus justo. Vestibulum risus lectus, adipiscing nec condimentum quis, condimentum nec nisl. Aliquam dictum sagittis velit sed iaculis. Morbi tristique augue sit amet nulla pulvinar id facilisis ligula mollis. Nam elit libero, tincidunt ut aliquam at, molestie in quam. Aenean rhoncus vehicula hendrerit.

\section{Thesis Outline}

The remainder of this thesis will be organized as follows :
\begin{description}
\item\textbf{Chapter~\ref{chapter:SemanticWebResources} - Semantic Web Resources:} presents different semantic web resources, their structures and how to browse them and how are they used in the literature. 
\item\textbf{Chapter~\ref{chapter:Measures} - Measures:} provides a solid background on semantic similarity and distance measures. We explore different metrics illustrating their pro/cons with examples.
\item\textbf{Chapter~\ref{chapter:Disambiguation} - Disambiguation:} reviews the literature and assesses the most relevant ways to disambiguate a list of keywords which may be organize into sentences or not. 
\item\textbf{Chapter~\ref{chapter:ExistingServices} - Existing Services:} describes existing image annotation services as well as their use-cases.
\item\textbf{Chapter~\ref{chapter:SotAConclusion} - State of the Art Conclusion:} summarizes the findings of the previous state of the art and opens the way to the presented contribution.
\item\textbf{Chapter~\ref{chapter:Methodology} - Proposed Methodology :} presents the chosen methodology as well as some organizational points.
\item\textbf{Chapter~\ref{chapter:Architecture} - Proposed Architecture :} details the technological choices by comparing them to their competitors and the chosen DBMS schema. Illustration figures will be presented.
\item\textbf{Chapter~\ref{chapter:Experiments} - Proposed Architecture :} presents the chosen dataset, details some of the main algorithms and reviews the tests' results with the use of different evaluation methods.
\item\textbf{Chapter~\ref{chapter:ContributionConclusion} - Contribution Conclusion :} summarize the findings of the presented research problem.
\end{description}
\part{State of the Art}
%!TEX root = ../thesis.tex

\chapter{Semantic Web} % Main chapter title

\label{chapter:SemanticWebResources} % Change X to a consecutive number; for referencing this chapter elsewhere, use \ref{ChapterX}

\lhead{Chapter 2 \emph{Semantic Web}} % Change X to a consecutive number; this is for the header on each page - perhaps a shortened title

Our study is focus on semantic enrichment of an initial set of keywords which can be organized as sentences or not. It is important to first understand what is a semantic concept and how concepts are organized into ontologies.\\ 

In this section, we will present some general notions about semantic concepts and review several semantic resources, their hierarchical structures and how we access them.

\section{Generalities} % (fold)
\label{sec:generalities}
Linguistic semantics is the study of meaning that is used for understanding human expression through language. It is easy for two human-being to communicate (given that they speak the same language) and to understand what their partner say even if he's using a tricky turn of phrase. However, this task becomes way more difficult when it comes to the comprehension of the human language by a machine. How can the computer guess that \say{I am totally dead} means in fact \say{I am really tired} and that the speaker isn't actually dead ? Machines need structured resources to understand us and the Semantic Web is one of them.\\

The notion of \say{Semantic Web} has been mentioned for the first time by Berners-Lee et al in \cite{berners2001semantic}. In this paper, they describe it as a Web which is readable by machines in opposite of most of Web's content which were designed for humans to read. The Semantic Web isn't a separate Web but an extension of the current one which will bring structure to the meaningful content of Web pages.\\

Two main technologies are used for the development of the Semantic Web : eXtensible Markup Language ( short XML) and the Resource Description Framework (short RDF). XML allows everyone to create their own tags and to arbitrary structure their documents but gives no information about what this structure means. Meaning is provided by RDF which stores it in sets of triples which are composed of a subject, a predicate and an object. Those three components can be related to the subject, the verb and object of an elementary sentence. In \cite{miller1998introduction}, Miller presents a short introduction to the RDF standard and precises that a \say{Resource} can be any object which is uniquely identifiable by a Uniform Resource Identifier (URI).\\

The third basic component of the Semantic Web are collections of information called ontologies. An ontology is, in computer science, a document which defines the relations among concepts. Basically, Web ontologies are composed of a taxonomy, which defines classes of objects and their relations, and a set of inference rules.\\

In addition, the The New Oxford Dictionary of English defines the notion of \say{semantic concept} as : \emph{An idea or thought that corresponds to some distinct entity or class of entities, or to its essential features, or determines the application of a term, and thus plays a part in the use of reason or language}.\\

Given those basic notions, we will now further detail four semantic web resources, their taxonomies and review some of their usage found in the literature.
% section generalities (end)

\section{Semantic Resources} % (fold)
\label{sec:semantic_resources}
\subsection{DBpedia} % (fold)
\label{sub:dbpedia}

DBpedia\footnote{http://wiki.dbpedia.org/} is a project originally launched by two German universities (Berlin and Leipzig) and backed by an important community. It explores Wikipedia\footnote{https://en.wikipedia.org/} and extract information from it which results on the creation of a multilingual, large-scale knowledge base. The extraction framework, all the available end-points as well as some facts and figures about the project are presented in \cite{lehmann2014dbpedia}.\\

DBpedia's ontology is based on classes (320 items) which form a subsumption hierarchy, the root element being \emph{owl:Thing}, with a maximal depth of 5\footnote{Complete classes tree : http://mappings.dbpedia.org/server/ontology/classes/}. Theses classes are described by a total of 1650 different properties, forming a large set of RDF triples (580 million extracted from the English version of Wikipedia).\\

Even though DBpedia is now a worldwide project and provides pages in 125 languages, the English one is still the most represented. We can indeed find 4.58 million of things\footnote{http://wiki.dbpedia.org/about/facts-figures} including 1,445,000 instances of the class \textit{Person}, 735,000 \textit{Place}, 251,000 \textit{Species} \dots The number of instances described in this language is about three time larger than the second and third language (French and German). \\

As well as any RDF-structured dataset, DBpedia can be requesting with SPARQL (which is an recursive acronym : SPARQL Protocol and RDF Query Language) queries. SPARQL allows the user to search, add, modify or delete RDF data available on the Internet, see \cite{prud2008sparql} for more details about the language. \\

\lstinputlisting[language=SPARQL,
	caption=SPARQL Query : Search classes,
	label={code:findClass}]
	{./Primitives/findClass.sparql}

\lstinputlisting[language=SPARQL,
	caption=SPARQL Query : Search superclasses,
	label={code:findSuperClass}]
	{./Primitives/findSuperClass.sparql}

Codes \ref{code:findClass} and \ref{code:findSuperClass} present two simple and generic SPARQL search queries which return respectively the class(es) of a given entity E and the superclass(es) of a given class C using the \textit{type} and \textit{subClassOf} predicates.\\

DBpedia also provides useful web services and HTTP endpoints. DBpedia Spotlight, which highlight DBpedia concepts in an input text is described in \cite{mendes2011dbpedia} and further details about disambiguation using this service are presented in subsection \ref{sec:dbpedia_spotlight}. The official DBpedia SPARQL endpoint\footnote{http://dbpedia.org/sparql} allows the user to send SPARQL queries to the online Virtuoso Triple Store by using the browser interface or by sending a HTPP request. We learn in \cite{lehmann2014dbpedia} that the average amount of hits per day of this endpoint is of 2,910,410 for the 3.8 dataset version. \\

We find lot of papers in the literature which mention DBpedia as an asset for systems based on the Semantic Web. If some of those papers are still in the research field, like \cite{passant2010dbrec} which proposes a music recommendation system built on top of DBpedia, others present \say{real} applications which are currently in use. We can cite for instance \cite{kobilarov2009media} which describes how the BBC\footnote{http://www.bbc.co.uk/} uses DBpedia to backbone its publications. 
% subsection dbpedia (end)

\subsection{Geonames} % (fold)
\label{sub:geonames}

GeoNames\footnote{http://www.geonames.org/} is a geographical database which contains over 10 million geographical names. The most documented countries\footnote{http://www.geonames.org/statistics/} are the United States of America (2,203,094 names), Norway (600,008) and China (526,456). All resources are categorized into one out of nine classes and further subcategorized into one out of 645 codes\footnote{http://www.geonames.org/export/codes.html}. Obviously, the root element of Geonames' hierarchy is mother earth.\\

Like DBpedia, the Geonames' ontology makes possible the addition of new resources to the World Wide Web. However, Geonames distinguish the feature's Concept from the RDF document about it. In consequence, each feature possesses two representation in Geonames. See the two following URIs as example :
\begin{itemize}
\item URI1 : http://sws.geonames.org/8015555/
\item URI2 : http://sws.geonames.org/8015555/about.rdf
\end{itemize}
The first one stands for the \say{Notre Dame de Fourvière} church in Lyon, France. This URI is used when one wants to refers to the church itself. The second URI is the RDF document with what Geonames knows about the church, its latitude and longitude for instance, or some nearby locations.\\

In order to allow the user to browse its \emph{Tremendous set of data} (Sir T. Berners-Lee), Geonames proposes a lot of REST-based web-services. In the case of image annotation, one in particular could be useful. Given that lot of recent numeric images contain EXIF data which embed GPS information such as the longitude and the latitude (see \cite{tevsic2005metadata} for further details about the EXIF format), one could add geographic tags to the picture with the help of the \textit{findNearbyPlaceName} service. Here is an example of usage, the response is display in Code \ref{code:passauXML}.\\
\underline{Query} : findNearbyPlaceName?lat=48.566\&lng=13.43\&username=demo 
\lstinputlisting[language=XML,
	caption=Geonames' XML response,
	float,
	label={code:passauXML}]
	{./Primitives/passau.xml}

Given its specific domain, we found less use-cases of Geonames in the literature than DBpedia's. Some interesting papers have however been presented, like \cite{stern2010detection} which use several semantic resources including Geonames in order to detect named entities in \say{Agence France Presse} (AFP) wires.
% subsection geonames (end)

\subsection{WordNet} % (fold)
\label{sub:wordnet}

WordNet\footnote{https://wordnet.princeton.edu/} is a lexical database of English which has been presented for the first time in 1995 in \cite{miller1995wordnet}. It is hosted by the Princeton University, currently running version 3.1 but there are no current plan for a future release due to limited staffing.\\

Its structure is based on the concept of \say{synset} (synonym set), a set cognitive synonyms. WordNet distinguish among Types (common nouns, verbs\dots) and Instances (specific persons\dots). Synsets are interlinked using conceptual, semantic and lexical relations. The hierarchy is built by the use of the super-subordinate relation (or hyperonymy, hyponymy in WordNet's jargon). These relations implements the two directions of the \say{IS-A} expression. For instance, \textit{fruit} is \textbf{hyperonym} of \textit{apple} and \textit{horse} is a \textbf{hyponym} of \textit{animal}, the root element being \say{entity}. Other relations are also provided, like the antonymy (opposite of synonymy) or the meronymy and its opposite holonymy which implements the \say{IS-PART-OF} relation : \textit{finger} is a \textbf{meronym} of \textit{hand}. All these relations are transitive.\\

This resource is useful if we are searching for entities. Since the maximal depth is of the ontology is of 16, the leafs are very detailed nouns (tsetse-fly, Yukon white birch, \dots) but it also contains more general concepts (vehicle, animal, \dots). WordNet contains at the moment 155,287 unique strings incuding 117,798 nouns.\\ 

It exists several ways to browse this resource. An online interface allows the user to manually query the dataset and to navigate in it through hyperlinks. For software and research purposes, the user has to download one of the released version of WordNet's dataset as well as a specific library according to the code language he's using.\\

Due to the lot of different relations between its synsets, WordNet has mostly been used for measuring semantic distances (see \cite{budanitsky2001semantic}, \cite{richardson1994using} and \cite{agirre2009study}) or to disambiguate texts (in \cite{resnik1995disambiguating}, \cite{voorhees1993using} and \cite{banerjee2002adapted}).\\
% subsection wordnet (end)

\subsection{ImageNet} % (fold)
\label{sub:imagenet}

ImageNet\footnote{http://www.image-net.org/} is an image database built on the WordNet hierarchy. It has been launch in 2009 and presented in \cite{deng2009imagenet}. Each of WordNet's synsets is depicted in ImageNet by a set of pictures (more than 500 in average). At the moment 14,197,122 images are referenced.\\

The stated aim of ImageNet is to serve as an asset for pedagogical and research purposes. It has for instance been used in \cite{deselaers2011visual} to measure the correlation between visual and semantic similarities or in a kindergarten in Canada to provide matching exercises to the children.\\

Despite the fact that this resource is not as used as the previous ones, it was important to cite it in order to highlight that semantic concepts and images can be related even tough they are different kinds of medium.  
% section imagenet (end)

% section semantic_resources (end)

\chapter{Measures} % Main chapter title

\label{Chapter4} % Change X to a consecutive number; for referencing this chapter elsewhere, use \ref{ChapterX}

\lhead{Chapter 4 \emph{Measures}} % Change X to a consecutive number; this is for the header on each page - perhaps a shortened title

%----------------------------------------------------------------------------------------
%	SECTION 1
%----------------------------------------------------------------------------------------

\section{Distance measures}

Lorem ipsum dolor sit amet, consectetur adipiscing elit. Aliquam ultricies lacinia euismod. Nam tempus risus in dolor rhoncus in interdum enim tincidunt. Donec vel nunc neque. In condimentum ullamcorper quam non consequat. Fusce sagittis tempor feugiat. Fusce magna erat, molestie eu convallis ut, tempus sed arcu. Quisque molestie, ante a tincidunt ullamcorper, sapien enim dignissim lacus, in semper nibh erat lobortis purus. Integer dapibus ligula ac risus convallis pellentesque.

%----------------------------------------------------------------------------------------
%	SECTION 2
%----------------------------------------------------------------------------------------

\section{Similarity measures}

Sed ullamcorper quam eu nisl interdum at interdum enim egestas. Aliquam placerat justo sed lectus lobortis ut porta nisl porttitor. Vestibulum mi dolor, lacinia molestie gravida at, tempus vitae ligula. Donec eget quam sapien, in viverra eros. Donec pellentesque justo a massa fringilla non vestibulum metus vestibulum. Vestibulum in orci quis felis tempor lacinia. Vivamus ornare ultrices facilisis. Ut hendrerit volutpat vulputate. Morbi condimentum venenatis augue, id porta ipsum vulputate in. Curabitur luctus tempus justo. Vestibulum risus lectus, adipiscing nec condimentum quis, condimentum nec nisl. Aliquam dictum sagittis velit sed iaculis. Morbi tristique augue sit amet nulla pulvinar id facilisis ligula mollis. Nam elit libero, tincidunt ut aliquam at, molestie in quam. Aenean rhoncus vehicula hendrerit.

%!TEX root = ../thesis.tex

\chapter{Disambiguation} % Main chapter title

\label{chapter:Disambiguation} % Change X to a consecutive number; for referencing this chapter elsewhere, use \ref{ChapterX}

\lhead{Chapter 3 \emph{Disambiguation}} % Change X to a consecutive number; this is for the header on each page - perhaps a shortened title

%----------------------------------------------------------------------------------------
%	SECTION 1
%----------------------------------------------------------------------------------------

\section{Section1}

Lorem ipsum dolor sit amet, consectetur adipiscing elit. Aliquam ultricies lacinia euismod. Nam tempus risus in dolor rhoncus in interdum enim tincidunt. Donec vel nunc neque. In condimentum ullamcorper quam non consequat. Fusce sagittis tempor feugiat. Fusce magna erat, molestie eu convallis ut, tempus sed arcu. Quisque molestie, ante a tincidunt ullamcorper, sapien enim dignissim lacus, in semper nibh erat lobortis purus. Integer dapibus ligula ac risus convallis pellentesque.
%!TEX root = ../thesis.tex

\chapter{Existing services} % Main chapter title

\label{chapter:ExistingServices} % Change X to a consecutive number; for referencing this chapter elsewhere, use \ref{ChapterX}

\lhead{Chapter 5 \emph{Existing services}} % Change X to a consecutive number; this is for the header on each page - perhaps a shortened title

We found in the literature lot of research which talk about multimedia tagging/annotation/labelling. These papers mainly use both the visual multimedia information (low-level features) and the textual information they might have at disposal (high-level features, meta-data). A review of the existing tagging-based applications is propose by Wand and al. in \cite{wang2012assistive}.\\
These works aren't our main focus, here we will present textual-only tagging approaches.

\section{Mixing Statistics and WordNet} % (fold)
\label{sec:mixing_statistics_and_wordnet}
In \cite{jin2005image}, Jin and al. propose the integration of the WordNet (\ref{sub:wordnet}) semantic resource in a statistic-based annotation process in order to remove irrelevant keywords.\\
Their algorithm is organized as follows : they first generate a set of keywords with the help of a statistical model called \emph{Translation Model} (TM). Some of those candidates tags are relevant and some aren't.\\
In order to filter the irrelevant ones, they then compute several semantic similarity measures : 
\begin{enumerate}
	\item Lin - \ref{sub:lin}
	\item Jiang and Conrath - \ref{sub:jiang_and_conrath}
	\item Banerjee and Pedersen - see \cite{banerjee2003extended}
\end{enumerate}
Finally, they combine these metrics using Dempster-Shafer Theory (\cite{shafer1976mathematical}) and the keywords with a resulting score under a chosen threshold are removed. Detailed method steps are presented in the original paper. \\
Concerning their results, they compare their \emph{TMHD} proposed approach with a basic TM process. Based on a set of most frequently used keywords, they found that, on average, precisions values of TM and TMHD are respectively 14.21\% and 33.11\%. This indicates that TMHD is 56.87\% better than TM. It is interesting to note that the recall score stays the same due to the fact that only irrelevant keywords are removed. They also compare \emph{TMHD} to the use of individual measures with TM and the results aren't as good as those from their combination.\\
These results show the power of knowledge-based data and similarity measures when it's added to a statistical model.
% section mixing_statistics_and_wordnet (end)

\section{Graph-cut based enrichment} % (fold)
\label{sec:graph_cut_based_enrichment}
In \cite{qian2011graph}, Qian and Hua expose their graph-based approach of the tag enrichment process. They represent each initial tag of their corpus as a node and interlink them (using \emph{n-links}). The weight of those n-links can be seen as the similarity between the two linked nodes, computed by the help of the Google distance \cite{cilibrasi2007google}. They add two virtual nodes called \emph{sink} and \emph{source}. Then, they link all nodes to one of these virtual nodes using \emph{t-links}.\\
The aim of their approach is to split all the tags into two distinct sets S (containing the source node) and T (containing the sink one) by assigning the labels s (source) if the tag is relevant to the image and t (sink) if not to the nodes. Then, they determine how many tags are relevant to the image by solving the combinational optimization problem through the graph.\\
This paper is really short and not very clear but it gives good ideas about the tags' representation as a graph and how to interlink them. According to the authors, the results are satisfactory.   
% section graph_cut_based_enrichment (end)

\section{Enrich Folksonomy Tag Space} % (fold)
\label{sec:enrich_folksonomy_tag_space}
Folksonomies are typical Web 2.0 systems that allow users to upload, tag and share content such as pictures, bookmarks \dots In \cite{angeletou2007bridging}, Angeletou and al. envisaged tag space enrichment with semantic relations by exploring online ontologies. Their method is composed of two phases :
\begin{itemize}
	\item Concept identification 
	\item Relation discovery
\end{itemize}
The first step is achieved by extracting concepts from online ontologies in which the local concept label matches the tag. In order to exploit all meanings, the authors retrieve all the potential semantic terms for each tag and then discover relation between them in the second phase. This means that no disambiguation is processed but it is a consequence of the relation discovery phase.\\
This phase consist of the identification of the relation between two tags \emph{T1} and \emph{T2}. Four kind of relations are distinguished :
\begin{itemize}
	\item Subsumption relation : \emph{T1} subClassOf \emph{T2}
	\item Disjointness relation : \emph{T1} disjointWith \emph{T2}
	\item Generic relation : Property1 hasDomain \emph{T1} and Property1 hasRange \emph{T2}
	\item Sibling relation : \emph{T1} and \emph{T2} share a common ancestor
	\item Instance Of relation : \emph{T1} instanceOf \emph{T2}
\end{itemize}
These relations can be found by two ways : a relation can be declared within an ontology or, if no ontology contain such relation, one is made by crossing knowledge from different ontologies.\\
The author the present different experiments as well as some issues rose during this phase. One in particular is important to keep in mind : when users tags resources, especially pictures, they tend to tag them with specific vocabulary, mainly instances rather than \emph{abstract} concepts. This can result on lot of \say{semantic noise} : tags which can't be match with concepts from online ontologies.\\
This paper is really interesting and approaches the topic in a very general point of view, which ensure the flexibility of its implementation. We will see in \ref{chapter:Methodology} that this method perfectly adapt to our study. 
% section enrich_folksonomy_tag_space (end)

%!TEX root = ../thesis.tex

\chapter{Conclusion} % Main chapter title

\label{chapter:SotAConclusion} % Change X to a consecutive number; for referencing this chapter elsewhere, use \ref{ChapterX}

\lhead{Chapter 7 \emph{Conclusion State of the Art}} % Change X to a consecutive number; this is for the header on each page - perhaps a shortened title

%----------------------------------------------------------------------------------------
%	SECTION 1
%----------------------------------------------------------------------------------------

\section{Section 1}

Lorem ipsum dolor sit amet, consectetur adipiscing elit. Aliquam ultricies lacinia euismod. Nam tempus risus in dolor rhoncus in interdum enim tincidunt. Donec vel nunc neque. In condimentum ullamcorper quam non consequat. Fusce sagittis tempor feugiat. Fusce magna erat, molestie eu convallis ut, tempus sed arcu. Quisque molestie, ante a tincidunt ullamcorper, sapien enim dignissim lacus, in semper nibh erat lobortis purus. Integer dapibus ligula ac risus convallis pellentesque.
\part{Contribution}
%!TEX root = ../thesis.tex

\chapter{Proposed Methodology} % Main chapter title

\label{chapter:Methodology} % Change X to a consecutive number; for referencing this chapter elsewhere, use \ref{ChapterX}

\lhead{Chapter 8 \emph{Proposed Methodology}} % Change X to a consecutive number; this is for the header on each page - perhaps a shortened title

Here we will present the chosen methodology, highly based on the three approaches presented in \ref{chapter:ExistingServices}. This methodology aims to support different kind of experiments in order to try to answer our thesis questions previously stated in \ref{sec:thesis_objectives}. 

\section{Global process} % (fold)
\label{sec:global_process}
With this work we want to propose a prototype which semantically enrich images given an initial set of tags. This prototype will be based on 3 steps :
\begin{enumerate}
	\item Concept identification
	\item Relation discovery
	\item Candidates detection
\end{enumerate}
As previously said, the two first steps will be similar as those presented in \ref{sec:enrich_folksonomy_tag_space}. \\

We also want to use a lemmatization process in order to group together the different inflected forms of a word so they can be analyzed as a single item. This will be achieve by using a hand-made tool based on regular expressions.\\

In order to compute semantic metrics, we will add virtual nodes to the graph, fulfiling the requirement of the Wu-Palmer evolved measure presented in \ref{ssub:evolved_wu_palmer}. More details about the graph structure will be presented in \ref{sec:graph_structure}.\\

We will set up different experiments, some based on the graph structure and some not. They will be run on a concrete images database and evaluated by several ways. Their results will then be presented, compared and discussed.
% section global_process (end)
\section{Knowledge bases} % (fold)
\label{sec:knowledge_bases}
The difference we want to propose between our tool and the approaches presented in \ref{chapter:ExistingServices} is to use several online ontologies in order to detect concepts and to create relations between them. We selected two resources : DBpedia and WordNet, already presented in \ref{sub:dbpedia} and \ref{sub:wordnet}. We will need ways to query these ontologies and browse them. In order not to get two separated graphs, we will create interlinks between concepts from both of the resources.\\
Eventually, the candidates we will propose as new annotations will also come from both of the ontologies. 
% section knowledge_bases (end)

\section{Data representation} % (fold)
\label{sec:data_representation}
As in \ref{sec:graph_cut_based_enrichment}, we want to represent our data as a graph. We will extract concepts from our chosen ontologies and save them as vertexes of our graph. We will interlink them using the already existing relations in their respective ontologies and add new ones we will detect by ourself. The graph structure is detailed in \ref{sec:graph_structure}.
% section data_representation (end)

%!TEX root = ../thesis.tex

\chapter{Proposed Architecture} % Main chapter title

\label{chapter:Architecture} % Change X to a consecutive number; for referencing this chapter elsewhere, use \ref{ChapterX}

\lhead{Chapter 9 \emph{Proposed Architecture}} % Change X to a consecutive number; this is for the header on each page - perhaps a shortened title

%----------------------------------------------------------------------------------------
%	SECTION 1
%----------------------------------------------------------------------------------------

\section{Technology Choices}

\subsection{Java}
Lorem ipsum dolor sit amet, consectetur adipiscing elit. Aliquam ultricies lacinia euismod. Nam tempus risus in dolor rhoncus in interdum enim tincidunt. Donec vel nunc neque. In condimentum ullamcorper quam non consequat. Fusce sagittis tempor feugiat. Fusce magna erat, molestie eu convallis ut, tempus sed arcu. Quisque molestie, ante a tincidunt ullamcorper, sapien enim dignissim lacus, in semper nibh erat lobortis purus. Integer dapibus ligula ac risus convallis pellentesque.

\subsection{Neo4j}
Lorem ipsum dolor sit amet, consectetur adipiscing elit. Aliquam ultricies lacinia euismod. Nam tempus risus in dolor rhoncus in interdum enim tincidunt. Donec vel nunc neque. In condimentum ullamcorper quam non consequat. Fusce sagittis tempor feugiat. Fusce magna erat, molestie eu convallis ut, tempus sed arcu. Quisque molestie, ante a tincidunt ullamcorper, sapien enim dignissim lacus, in semper nibh erat lobortis purus. Integer dapibus ligula ac risus convallis pellentesque.

\subsection{NLP}
Lorem ipsum dolor sit amet, consectetur adipiscing elit. Aliquam ultricies lacinia euismod. Nam tempus risus in dolor rhoncus in interdum enim tincidunt. Donec vel nunc neque. In condimentum ullamcorper quam non consequat. Fusce sagittis tempor feugiat. Fusce magna erat, molestie eu convallis ut, tempus sed arcu. Quisque molestie, ante a tincidunt ullamcorper, sapien enim dignissim lacus, in semper nibh erat lobortis purus. Integer dapibus ligula ac risus convallis pellentesque.

\subsection{DBpedia Spotlight}
Lorem ipsum dolor sit amet, consectetur adipiscing elit. Aliquam ultricies lacinia euismod. Nam tempus risus in dolor rhoncus in interdum enim tincidunt. Donec vel nunc neque. In condimentum ullamcorper quam non consequat. Fusce sagittis tempor feugiat. Fusce magna erat, molestie eu convallis ut, tempus sed arcu. Quisque molestie, ante a tincidunt ullamcorper, sapien enim dignissim lacus, in semper nibh erat lobortis purus. Integer dapibus ligula ac risus convallis pellentesque.

\subsection{JAWS}
Lorem ipsum dolor sit amet, consectetur adipiscing elit. Aliquam ultricies lacinia euismod. Nam tempus risus in dolor rhoncus in interdum enim tincidunt. Donec vel nunc neque. In condimentum ullamcorper quam non consequat. Fusce sagittis tempor feugiat. Fusce magna erat, molestie eu convallis ut, tempus sed arcu. Quisque molestie, ante a tincidunt ullamcorper, sapien enim dignissim lacus, in semper nibh erat lobortis purus. Integer dapibus ligula ac risus convallis pellentesque.

\subsection{JENA}
Lorem ipsum dolor sit amet, consectetur adipiscing elit. Aliquam ultricies lacinia euismod. Nam tempus risus in dolor rhoncus in interdum enim tincidunt. Donec vel nunc neque. In condimentum ullamcorper quam non consequat. Fusce sagittis tempor feugiat. Fusce magna erat, molestie eu convallis ut, tempus sed arcu. Quisque molestie, ante a tincidunt ullamcorper, sapien enim dignissim lacus, in semper nibh erat lobortis purus. Integer dapibus ligula ac risus convallis pellentesque.

\subsection{JSoup}
Lorem ipsum dolor sit amet, consectetur adipiscing elit. Aliquam ultricies lacinia euismod. Nam tempus risus in dolor rhoncus in interdum enim tincidunt. Donec vel nunc neque. In condimentum ullamcorper quam non consequat. Fusce sagittis tempor feugiat. Fusce magna erat, molestie eu convallis ut, tempus sed arcu. Quisque molestie, ante a tincidunt ullamcorper, sapien enim dignissim lacus, in semper nibh erat lobortis purus. Integer dapibus ligula ac risus convallis pellentesque.

%----------------------------------------------------------------------------------------
%	SECTION 2
%----------------------------------------------------------------------------------------

\section{Graph Structure}

\subsection{Pro-Cons}
Lorem ipsum dolor sit amet, consectetur adipiscing elit. Aliquam ultricies lacinia euismod. Nam tempus risus in dolor rhoncus in interdum enim tincidunt. Donec vel nunc neque. In condimentum ullamcorper quam non consequat. Fusce sagittis tempor feugiat. Fusce magna erat, molestie eu convallis ut, tempus sed arcu. Quisque molestie, ante a tincidunt ullamcorper, sapien enim dignissim lacus, in semper nibh erat lobortis purus. Integer dapibus ligula ac risus convallis pellentesque.

\subsection{Vertexes}
Lorem ipsum dolor sit amet, consectetur adipiscing elit. Aliquam ultricies lacinia euismod. Nam tempus risus in dolor rhoncus in interdum enim tincidunt. Donec vel nunc neque. In condimentum ullamcorper quam non consequat. Fusce sagittis tempor feugiat. Fusce magna erat, molestie eu convallis ut, tempus sed arcu. Quisque molestie, ante a tincidunt ullamcorper, sapien enim dignissim lacus, in semper nibh erat lobortis purus. Integer dapibus ligula ac risus convallis pellentesque.

\subsection{Edges}
Lorem ipsum dolor sit amet, consectetur adipiscing elit. Aliquam ultricies lacinia euismod. Nam tempus risus in dolor rhoncus in interdum enim tincidunt. Donec vel nunc neque. In condimentum ullamcorper quam non consequat. Fusce sagittis tempor feugiat. Fusce magna erat, molestie eu convallis ut, tempus sed arcu. Quisque molestie, ante a tincidunt ullamcorper, sapien enim dignissim lacus, in semper nibh erat lobortis purus. Integer dapibus ligula ac risus convallis pellentesque.
\chapter{Experiments} % Main chapter title

\label{Chapter10} % Change X to a consecutive number; for referencing this chapter elsewhere, use \ref{ChapterX}

\lhead{Chapter 10 \emph{Experiments}} % Change X to a consecutive number; this is for the header on each page - perhaps a shortened title

%----------------------------------------------------------------------------------------
%	SECTION 1
%----------------------------------------------------------------------------------------

\section{Dataset}

Lorem ipsum dolor sit amet, consectetur adipiscing elit. Aliquam ultricies lacinia euismod. Nam tempus risus in dolor rhoncus in interdum enim tincidunt. Donec vel nunc neque. In condimentum ullamcorper quam non consequat. Fusce sagittis tempor feugiat. Fusce magna erat, molestie eu convallis ut, tempus sed arcu. Quisque molestie, ante a tincidunt ullamcorper, sapien enim dignissim lacus, in semper nibh erat lobortis purus. Integer dapibus ligula ac risus convallis pellentesque.

%----------------------------------------------------------------------------------------
%	SECTION 2
%----------------------------------------------------------------------------------------

\section{Code Explanation}

Sed ullamcorper quam eu nisl interdum at interdum enim egestas. Aliquam placerat justo sed lectus lobortis ut porta nisl porttitor. Vestibulum mi dolor, lacinia molestie gravida at, tempus vitae ligula. Donec eget quam sapien, in viverra eros. Donec pellentesque justo a massa fringilla non vestibulum metus vestibulum. Vestibulum in orci quis felis tempor lacinia. Vivamus ornare ultrices facilisis. Ut hendrerit volutpat vulputate. Morbi condimentum venenatis augue, id porta ipsum vulputate in. Curabitur luctus tempus justo. Vestibulum risus lectus, adipiscing nec condimentum quis, condimentum nec nisl. Aliquam dictum sagittis velit sed iaculis. Morbi tristique augue sit amet nulla pulvinar id facilisis ligula mollis. Nam elit libero, tincidunt ut aliquam at, molestie in quam. Aenean rhoncus vehicula hendrerit.

\section{Results and Analysis}

\subsection{Evaluation methodology}
Lorem ipsum dolor sit amet, consectetur adipiscing elit. Aliquam ultricies lacinia euismod. Nam tempus risus in dolor rhoncus in interdum enim tincidunt. Donec vel nunc neque. In condimentum ullamcorper quam non consequat. Fusce sagittis tempor feugiat. Fusce magna erat, molestie eu convallis ut, tempus sed arcu. Quisque molestie, ante a tincidunt ullamcorper, sapien enim dignissim lacus, in semper nibh erat lobortis purus. Integer dapibus ligula ac risus convallis pellentesque.

\subsection{Graph-based experiments}
\subsubsection{Direct Neighbors}
\subsubsection{Lists - WL}
\subsubsection{Lists - SL}

\subsection{Plain-text experiments}
\subsubsection{WikiLinks}
\subsubsection{WikiContent}

\chapter{Conclusion} % Main chapter title

\label{Chapter11} % Change X to a consecutive number; for referencing this chapter elsewhere, use \ref{ChapterX}

\lhead{Chapter 11 \emph{Contribution Conclusion}} % Change X to a consecutive number; this is for the header on each page - perhaps a shortened title

%----------------------------------------------------------------------------------------
%	SECTION 1
%----------------------------------------------------------------------------------------

\section{SEction 1}

Lorem ipsum dolor sit amet, consectetur adipiscing elit. Aliquam ultricies lacinia euismod. Nam tempus risus in dolor rhoncus in interdum enim tincidunt. Donec vel nunc neque. In condimentum ullamcorper quam non consequat. Fusce sagittis tempor feugiat. Fusce magna erat, molestie eu convallis ut, tempus sed arcu. Quisque molestie, ante a tincidunt ullamcorper, sapien enim dignissim lacus, in semper nibh erat lobortis purus. Integer dapibus ligula ac risus convallis pellentesque.


%----------------------------------------------------------------------------------------
%	THESIS CONTENT - APPENDICES
%----------------------------------------------------------------------------------------

\addtocontents{toc}{\vspace{2em}} % Add a gap in the Contents, for aesthetics

\appendix % Cue to tell LaTeX that the following 'chapters' are Appendices

% Include the appendices of the thesis as separate files from the Appendices folder
% Uncomment the lines as you write the Appendices

% Appendix A

\chapter{Appendix Title Here} % Main appendix title

\label{AppendixA} % For referencing this appendix elsewhere, use \ref{AppendixA}

\lhead{Appendix A. \emph{Appendix Title Here}} % This is for the header on each page - perhaps a shortened title

Write your Appendix content here.
%\input{Appendices/AppendixB}
%\input{Appendices/AppendixC}

\addtocontents{toc}{\vspace{2em}} % Add a gap in the Contents, for aesthetics

\backmatter

%----------------------------------------------------------------------------------------
%	BIBLIOGRAPHY
%----------------------------------------------------------------------------------------

\label{Bibliography}

\lhead{\emph{Bibliography}} % Change the page header to say "Bibliography"

\bibliographystyle{unsrtnat} % Use the "unsrtnat" BibTeX style for formatting the Bibliography

\bibliography{Bibliography} % The references (bibliography) information are stored in the file named "Bibliography.bib"

\end{document}  